خواسته‌ی کارمی درست است: در تیمی که چند نفر هم‌زمان روی منو، سفارش و هماهنگی آشپزخانه کار می‌کنند، تنها جایی که می‌توان جلوی ورود مشکل به شاخه‌ی اصلی را گرفت، {قبل از ادغام} است. بعد از ادغام، هزینه‌ی پیدا کردن مقصر و برگرداندن تغییر چند برابر می‌شود.

تفاوت پایپلاین \lr{merge request} با پایپلاین شاخه‌ی اصلی:

\begin{center}
\begin{tabular}{|l|l|l|}
\hline
 & {پایپلاین \lr{MR}} & {پایپلاین \lr{main}} \\
\hline
هدف & آیا اجازه‌ی ورود دارد؟ & آیا آماده‌ی انتشار است؟ \\
\hline
کد مورد آزمون & نتیجه‌ی {ادغام} شاخه با \lr{main} & خود \lr{main} \\
\hline
استقرار & هیچ (یا حداکثر \lr{review app} موقت) & \lr{staging} و سپس \lr{production} \\
\hline
انتشار \lr{artifact} & فقط بیلد، بدون پوش پایدار & پوش ایمیج با تگ نسخه \\
\hline
دامنه‌ی تست & تست‌های سریع + تست‌های متأثر & کل مجموعه، شامل \lr{e2e} کامل \\
\hline
سیاست شکست & جلوی \lr{merge} را می‌گیرد & هشدار فوری به تیم، مسدودسازی انتشار \\
\hline
\end{tabular}
\end{center}

مهم‌ترین تفاوت فنی، همان سطر دوم است. اجرای پایپلاین روی آخرین کامیت شاخه‌ی توسعه‌دهنده کافی نیست، چون ممکن است آن شاخه سبز باشد ولی بعد از ادغام با \lr{main} خراب شود (\lr{semantic conflict}: مثلا یک نفر امضای تابع محاسبه‌ی قیمت سفارش را عوض کرده و نفر دیگر هم‌زمان یک فراخوانی جدید به همان تابع اضافه کرده است؛ گیت تعارض نمی‌بیند ولی کد کامپایل نمی‌شود). بنابراین پایپلاین باید روی \lr{merge result} اجرا شود و اگر \lr{main} جلو رفت، دوباره اجرا شود.

 مراحل پایپلاین \lr{merge request}:

\begin{enumerate}
	\item {بررسی‌های ارزان}: قالب پیام کامیت، نبود فایل حجیم، \lr{secret detection}، لینت فقط روی فایل‌های تغییرکرده. این‌ها اول اجرا می‌شوند تا بازخورد تقریبا بلافاصله برسد.
	\item {\lr{lint} و بررسی تایپ}: \lr{formatter}، لینتر و کامپایل تایپ روی کل پروژه.
	\item {\lr{unit test}}: سریع، بدون وابستگی خارجی، به همراه گزارش پوشش و مقایسه با \lr{main} (کاهش پوشش، خودش یک هشدار است).
	\item {\lr{build}}: بیلد ایمیج سرویس‌ها. اگر کد کامپایل شود ولی ایمیج ساخته نشود، مشکل واقعی است و باید همین‌جا دیده شود.
	\item {\lr{integration test}}: با بالا آوردن وابستگی‌های واقعی (پایگاه داده، \lr{Redis}، صف سفارش‌ها) به‌صورت سرویس کنار جاب؛ مسیرهای حساسی مثل ثبت سفارش و اعلان به آشپزخانه.
	\item {بررسی امنیتی}: \lr{SAST} روی \lr{diff}، اسکن وابستگی‌ها و اسکن ایمیج ساخته‌شده.
	\item {\lr{review app}}: یک نمونه‌ی موقت از کل سامانه با \lr{URL} یکتا برای هر \lr{MR}، تا کارمی بتواند تغییر منو را واقعا ببیند نه اینکه تصورش کند. با بسته شدن \lr{MR} خودکار حذف می‌شود.
	البته این مورد اختیاری هست و معمولاً استفاده نمی‌شود.
\end{enumerate}

{مسدودکننده (\lr{blocking}):}
\begin{itemize}
	\item شکست کامپایل یا بیلد ایمیج
	\item حداقل یکی از تست‌های واحد یا یکپارچگی شکست‌خورده باشد
	\item خطای لینت و \lr{formatter}
	\item آسیب‌پذیری \lr{high} یا \lr{critical} در وابستگی‌ها، و هر \lr{secret} پیدا شده در کد
	\item افت پوشش تست زیر آستانه‌ی توافق‌شده
\end{itemize}

{غیرمسدودکننده (فقط هشدار، \lr{allow\_failure}):}
\begin{itemize}
	\item آسیب‌پذیری \lr{low/medium} بدون وصله‌ی موجود
	\item هشدارهای کارایی مثل افزایش حجم \lr{bundle} یا زمان بیلد
\end{itemize}

معیار ساده: چیزی که {قطعیت} دارد و {قابل رفع توسط نویسنده‌ی همین تغییر} است، مسدودکننده باشد؛ چیزی که قضاوتی است یا رفعش به تصمیم تیم نیاز دارد، هشدار باشد.

 رساندن نتیجه به بازبین‌ها و کارمی:
\begin{itemize}
	\item {وضعیت روی خود \lr{MR}}: هر جاب یک \lr{status check} جداگانه با نام گویا؛ بازبین در یک نگاه می‌فهمد کدام مرحله خراب است، نه اینکه فقط قرمز ببیند.
	\item {قوانین حفاظت از شاخه}: دکمه‌ی \lr{merge} تا سبز شدن بررسی‌های اجباری غیرفعال باشد. اتکا به یادت باشد قبل از \lr{merge} نگاه کنی جواب نمی‌دهد.
	\item {گزارش‌های خلاصه به‌صورت کامنت خودکار}: تست‌های شکست‌خورده با نام و پیام، اختلاف پوشش نسبت به \lr{main}، و فهرست آسیب‌پذیری‌های تازه‌واردشده همه در قالب \lr{diff}، نه فهرست کامل؛ چون بازبین فقط مسئول {تغییر} است.
	\item {\lr{artifact}های قابل دانلود}: گزارش پوشش، خروجی تست‌ها، لاگ کامل هر جاب و لاگ تست‌های \lr{e2e} شکست‌خورده.
	\item {لینک \lr{review app}} مستقیم در توضیحات \lr{MR} برای بررسی دستی توسط کارمی.
	\item {قوانین تأیید}: حداقل یک تأیید از صاحب آن بخش (\lr{CODEOWNERS})؛ برای مثال تغییر در ماژول آشپزخانه نیاز به تأیید مسئول همان ماژول دارد.
\end{itemize}

با این ساختار، تصمیم کارمی مبتنی بر شواهد است: پایپلاین می‌گوید چه چیزی خراب است، بازبین می‌گوید آیا این راه‌حل درستی است، و دکمه‌ی \lr{merge} تنها وقتی فعال می‌شود که هر دو پاسخ مثبت باشند.
